% Conversão de ages-iv-sprint-4; origem imutável RELATORIO_ATUAL.odt.
\hypertarget{__RefHeading___Toc4853_3118176633}{}
\refstepcounter{subsection}
\subsection*{5.3.5 Sprint 4}
\label{ages-iv:text-h-0121}
\addcontentsline{toc}{subsection}{5.3.5 Sprint 4}

% user-ages-iv-sprint-4-report-p01
\indent A Sprint 4 marcou a etapa final de desenvolvimento do Recorda e foi, na minha percepção, o momento em que a equipe atingiu seu maior nível de organização e autonomia. Iniciei a Sprint realizando a decomposição das User Stories do Épico 7, último épico previsto para o projeto, em tarefas de frontend e backend. Fiquei responsável por especificar detalhadamente cada atividade, identificar suas dependências, configurar os bloqueios e organizar tudo no quadro Kanban do GitHub. Como havíamos antecipado boa parte do Épico 6 na Sprint anterior, conseguimos estruturar um cronograma que reservasse a última semana do projeto para testes, correções e refinamentos da aplicação.\par

% user-ages-iv-sprint-4-report-p02
\indent Também acompanhei os AGES III nas definições de infraestrutura, auxiliando na elaboração do orçamento para implantação dos serviços na AWS. Compartilhei experiências do semestre anterior, quando atuei como AGES III, principalmente sobre os desafios de conciliar decisões arquiteturais com os recursos financeiros limitados disponibilizados pela AGES. Paralelamente, pesquisei aspectos relacionados à publicação do aplicativo na App Store, como políticas de privacidade, termos de uso, exclusão de contas e disponibilização de credenciais para avaliação. Como não possuía experiência aprofundada nesse processo, conversei com o Andrei, AGES IV e integrante da Apple Developer Academy, para compreender melhor esses requisitos e possíveis impedimentos à publicação.\par

% user-ages-iv-sprint-4-report-p03
\indent No desenvolvimento, a Sprint foi especialmente tranquila. Diferentemente das primeiras entregas, praticamente não enfrentamos dificuldades relacionadas à compreensão das tarefas, à qualidade do código ou à integração entre funcionalidades. Os integrantes já conheciam os processos, tinham clareza sobre suas responsabilidades e conseguiam resolver os ajustes necessários com bastante independência. Essa fluidez foi resultado de um trabalho contínuo dos AGES IV na organização das squads, na especificação das atividades e no aperfeiçoamento dos processos ao longo do semestre.\par

% user-ages-iv-sprint-4-report-p04
\indent Uma das principais decisões de gestão foi antecipar o \textit{code freeze} das funcionalidades em uma semana, buscando garantir tempo suficiente para a implementação e execução dos testes unitários, de integração e de sistema. Isso tornou a conclusão das últimas funcionalidades mais corrida, mas permitiu que a equipe direcionasse seus esforços finais à qualidade e à estabilidade da aplicação. Como os dois AGES III estavam concentrados na configuração e implantação dos serviços da AWS, nós, AGES IV, assumimos grande parte das revisões de código. Nesse período, fui responsável pela revisão de nove Pull Requests, envolvendo funcionalidades tanto de frontend quanto de backend.\par

% user-ages-iv-sprint-4-report-p05
\indent Com as funcionalidades concluídas e integradas, concentramos a última semana no desenvolvimento de testes End-to-End, além de pequenos ajustes de desempenho, correções de bugs e melhorias de UX/UI. Em conjunto com os demais AGES IV, alinhamos que as atividades de testes e refinamento poderiam continuar até a entrega da Sprint 4, aproveitando esse período adicional sem comprometer a consolidação das funcionalidades.\par

% user-ages-iv-sprint-4-report-p06
\indent Também participei ativamente da concepção e elaboração dos slides para a apresentação aos stakeholders e para a apresentação final de todos os projetos da AGES. Planejamos cuidadosamente os roteiros das demonstrações, considerando especialmente a limitação de dez minutos para explicar o desenvolvimento do projeto e apresentar as principais funcionalidades do Recorda. Como a aplicação havia crescido consideravelmente ao longo do semestre, selecionar o que demonstrar e organizar uma apresentação clara dentro desse tempo exigiu bastante atenção e preparação.\par

% user-ages-iv-sprint-4-report-p07
\indent A apresentação final aos clientes foi um encerramento muito positivo para o trabalho desenvolvido durante o semestre. Conseguimos demonstrar a solução de maneira abrangente, contemplando as funcionalidades construídas ao longo das Sprints, e recebemos um retorno muito satisfatório dos stakeholders. Um dos momentos mais marcantes foi o feedback da professora orientadora Alessandra Dutra, que expressou orgulho pelo resultado alcançado e destacou que havíamos construído uma solução completa, e não apenas uma prova de conceito. Esse reconhecimento representou muito bem o amadurecimento do projeto e o esforço coletivo da equipe, que conseguiu transformar uma proposta inicial em uma aplicação com funcionalidades integradas e uma experiência de uso consistente.\par

\begin{reflection}
% user-ages-iv-sprint-4-reflection-p01
\noindent A Sprint 4 me fez perceber o quanto as retrospectivas e a melhoria contínua realmente fizeram diferença no desenvolvimento do projeto. Foi impressionante observar a sintonia da equipe e a naturalidade com que as atividades eram executadas, principalmente quando comparadas às dificuldades que enfrentávamos no início do semestre. Olhando para trás, percebo que essa evolução não aconteceu ao acaso. Cada problema discutido, cada ajuste na organização das squads, nas especificações das tarefas e nos processos de desenvolvimento contribuiu para chegarmos a esse resultado. Foi especialmente gratificante perceber que os \textit{action items} definidos nas retrospectivas não ficaram apenas no planejamento, mas produziram mudanças concretas na forma como o time trabalhava.\par

% user-ages-iv-sprint-4-reflection-p02
\noindent Encerrar minha participação como AGES IV com uma entrega tão positiva também trouxe um sentimento de realização. Ao longo do semestre, compreendi cada vez melhor que liderar uma equipe não significa necessariamente desenvolver mais do que os outros, mas criar condições para que todos consigam trabalhar bem. Ver o time alcançar esse nível de autonomia e receber o reconhecimento dos clientes e da professora foi uma confirmação muito significativa do valor desse trabalho de gestão, que nem sempre é tão visível quanto o código desenvolvido, mas que foi fundamental para o resultado que alcançamos.\par

\end{reflection}
